%==============================================================================
\section{Thuật toán học tăng cường sâu (Deep RL Algorithms)}
\label{sec:deep-rl-algorithms}
%==============================================================================

\begin{dinhnghia}[title={Thuật toán Deep RL}]
Thuật toán học tăng cường sâu là loại thuật toán học máy \textbf{kết hợp học sâu và học tăng cường}.
Deep RL giải quyết bài toán cho agent tính toán học ra quyết định bằng cách tích hợp học sâu từ dữ liệu đầu vào không có cấu trúc, \textbf{không cần} kỹ thuật thủ công thiết kế không gian trạng thái.
Các thuật toán này có khả năng quyết định hành động tối ưu target ngay cả khi đầu vào có kích thước lớn.
\end{dinhnghia}

\subsection{Học tăng cường (Reinforcement Learning)}

\begin{ynghia}[title={Vai trò của RL trong Deep RL}]
\textbf{Học tăng cường} gồm một agent học từ \textbf{phản hồi} đối với hành động của mình khi khám phá môi trường.
target chính của agent là \textbf{tối đa hóa phần thưởng tích lũy} bằng cách phát triển chiến lược định hướng ra quyết định trong mọi tình huống có thể.
\end{ynghia}

\subsection{Vai trò của học sâu trong RL}

\begin{luuy}[title={Hạn chế xấp xỉ truyền thống}]
Trong thuật toán RL truyền thống, \textbf{bảng tra} hoặc xấp xỉ hàm cơ bản thường được dùng để biểu diễn hàm giá trị, chính sách hoặc mô hình.
Các cách này \textbf{không đủ hiệu quả} trong bối cảnh khó như game video, robot hoặc xử lý ngôn ngữ tự nhiên.
\end{luuy}

\begin{phuongphap}[title={Nền tảng Deep RL}]
\textbf{Mạng nơ-ron} cho phép xấp xỉ các hàm \textbf{phức tạp, đa chiều} thông qua học sâu.
Đây là nền tảng của học tăng cường sâu.
\end{phuongphap}

\subsection{Lợi ích kết hợp học sâu và RL}

\begin{tinhchat}[title={Ba lợi ích chính}]
Kết hợp mạng học sâu với học tăng cường mang lại các lợi ích sau:
\begin{enumerate}
  \item Xử lý đầu vào \textbf{chiều cao} (ảnh thô, dữ liệu cảm biến liên tục).
  \item Hiểu quan hệ \textbf{phức tạp} giữa trạng thái và hành động thông qua học.
  \item Học \textbf{biểu diễn chung} bằng cách khái quát hóa giữa các trạng thái và hành động khác nhau.
\end{enumerate}
\end{tinhchat}

\subsection{Các thuật toán Deep RL phổ biến}

\begin{ynghia}[title={Tổng quan}]
Dưới đây là một số thuật toán học tăng cường sâu phổ biến.
\end{ynghia}

\subsubsection{Deep Q-Networks (DQN)}

\begin{dinhnghia}[title={Deep Q-Network}]
\textbf{Deep Q-Network (DQN)} là mở rộng của Q-learning cổ điển, dùng \textbf{mạng nơ-ron sâu} ước lượng hàm giá trị hành động $Q(s,a)$.
Thay vì lưu giá trị $Q$ trong bảng, DQN dùng mạng nơ-ron xử lý miền đầu vào phức tạp như \textbf{pixel game}.
Nhờ đó RL có thể giải quyết tác vụ phức tạp---ví dụ chơi Atari, agent học từ đầu vào \textbf{trực quan}.
\end{dinhnghia}

\begin{phuongphap}[title={Ổn định huấn luyện DQN}]
DQN cải thiện ổn định huấn luyện qua hai phương pháp chính:
\begin{itemize}
  \item \textbf{Experience replay}: lưu và lấy mẫu kinh nghiệm quá khứ;
  \item \textbf{Target network}: duy trì target $Q$ nhất quán bằng cách cập nhật định kỳ một mạng riêng.
\end{itemize}
Các cải tiến này giúp DQN học hiệu quả trong bối cảnh \textbf{quy mô lớn}.
\end{phuongphap}

\subsubsection{Double Deep Q-Networks (DDQN)}

\begin{ynghia}[title={Giảm thiên lệch overestimation}]
\textbf{Double Deep Q-Network (DDQN)} cải thiện DQN bằng cách giảm vấn đề \textbf{thiên lệch overestimation} khi cập nhật giá trị $Q$.
Trong DQN thông thường, một mạng $Q$ duy nhất vừa \textbf{chọn hành động} vừa \textbf{ước lượng giá trị}, dễ dẫn tới xấp xỉ giá trị \textbf{lạc quan quá mức}.
\end{ynghia}

\begin{phuongphap}[title={Hai mạng trong DDQN}]
DDQN dùng \textbf{hai mạng riêng}:
\begin{itemize}
  \item Mạng $Q$ \textbf{hiện tại} để chọn hành động;
  \item Mạng \textbf{target} để đánh giá hành động đó.
\end{itemize}
Việc tách vai trò giảm thiên lệch, nâng \textbf{độ chính xác} học.
DDQN vẫn kế thừa experience replay và target network từ DQN để tăng \textbf{độ bền} và \textbf{tin cậy}.
\end{phuongphap}

\subsubsection{Dueling Deep Q-Networks (Dueling DQN)}

\begin{dinhnghia}[title={Dueling DQN}]
\textbf{Dueling Deep Q-Networks} mở rộng DQN chuẩn bằng cách \textbf{tách} giá trị $Q$ thành hai thành phần:
\begin{itemize}
  \item Hàm giá trị trạng thái $V(s)$;
  \item Hàm \textbf{advantage} $A(s,a)$, ước lượng mức giá trị của từng hành động so với giá trị trung bình.
\end{itemize}
Giá trị $Q$ cuối cùng được ước lượng bằng cách \textbf{kết hợp} hai thành phần này.
\end{dinhnghia}

\begin{tinhchat}[title={Lợi ích biểu diễn tách rời}]
Biểu diễn này tăng \textbf{sức mạnh và hiệu quả} của Q-learning: mô hình ước lượng $V(s)$ chính xác hơn, đồng thời giảm nhu cầu ước lượng giá trị hành động chi tiết trong một số tình huống.
\end{tinhchat}

\subsubsection{Policy Gradient Methods}

\begin{dinhnghia}[title={Phương pháp Policy Gradient}]
\textbf{Policy Gradient} là nhóm thuật toán theo hướng \textbf{lặp chính sách}: điều chỉnh trực tiếp policy để đạt policy tối ưu \textbf{tối đa hóa phần thưởng kỳ vọng}.
Thay vì chỉ học hàm giá trị, các phương pháp này \textbf{tối ưu policy} theo gradient của target định nghĩa đối với tham số policy.
\end{dinhnghia}

\begin{phuongphap}[title={target và các biến thể}]
target chính là tính \textbf{gradient phần thưởng trung bình} và cập nhật chiến lược.
Các thuật toán tiêu biểu: \textbf{REINFORCE}, \textbf{Actor--Critic}, \textbf{Proximal Policy Optimization (PPO)}.
Các cách tiếp cận này hiệu quả trong không gian hành động \textbf{chiều cao} hoặc \textbf{liên tục}.
\end{phuongphap}

\subsubsection{Proximal Policy Optimization (PPO)}

\begin{dinhnghia}[title={PPO}]
\textbf{Proximal Policy Optimization (PPO)} là thuật toán RL nhằm tối ưu chính sách \textbf{ổn định và hiệu quả} hơn.
PPO cập nhật policy bằng cách tối đa hóa hàm target gắn với policy, nhưng \textbf{giới hạn} mức thay đổi policy cho phép để tránh biến động đột ngột.
\end{dinhnghia}

\begin{phuongphap}[title={target clipped}]
Policy mới \textbf{không được} lệch quá xa policy cũ; PPO dùng \textbf{target clipped} để đảm bảo điều đó.
Nhờ target clipped, PPO ngăn thay đổi policy quá lớn giữa hai lần cập nhật liên tiếp.
Sự cân bằng giữa khám phá và khai thác giúp tránh \textbf{suy giảm hiệu năng} và thúc đẩy \textbf{hội tụ mượt} hơn.
\end{phuongphap}

\begin{vidu}[title={Phạm vi ứng dụng PPO}]
PPO được áp dụng trong Deep RL cho cả không gian hành động \textbf{liên tục} và \textbf{rời rạc} nhờ tính \textbf{đơn giản} và \textbf{hiệu quả}.
\end{vidu}
